\subsection{Simulation-Based Relation Label Construction}
\label{app:relation_label_construction}

\paragraph{Simulator supervision and task-specific assets.}
At each environment step, we construct the scene graph from privileged
simulator state. We use object poses, the end-effector pose, gripper width,
contact forces, and the simulator grasp predicate only to construct graph
labels; these quantities are not provided directly to the world model or the
policy. Before training, the task-mining phase processes successful rollouts
once to obtain the task-relevant object set, the admissible interaction types
for each entity, object-local affordance components, and relation-specific
discretization thresholds. These assets are then fixed throughout training.

An admissible relation must have a valid endpoint type, be supported by the
mined interaction types of its endpoints, and have the required affordance
components and calibration thresholds. A relation is emitted only when every
object endpoint lies within a camera view. The end effector need not be
visually observed because its state is available from robot proprioception.
When an object leaves the camera views, we retain its vertex with zero bounding
boxes but omit its incident relations. We do not carry forward their previous
labels; if the relation later becomes admissible again, it is recomputed from
the current state.

\paragraph{Relation vocabulary.}
Table~\ref{tab:main_relation_vocab} summarizes the relation vocabulary.
Physical-state relations describe realized interactions. Spatial relations
describe coarse pairwise geometry, whereas affordance relations measure
whether the current relative configuration is suitable for an interaction.
All relations carry an absolute state, and spatial and affordance relations
additionally carry a temporal-change state.

\begin{table}[h]
\centering
\footnotesize
\setlength{\tabcolsep}{4pt}
\renewcommand{\arraystretch}{1.08}
\caption{Relation vocabulary and admissible endpoint types. Spatial and
affordance relations carry both absolute and temporal-change states; physical
relations carry only an absolute state.}
\label{tab:main_relation_vocab}
\begin{tabularx}{0.99\linewidth}{@{}
>{\raggedright\arraybackslash}p{0.18\linewidth}
>{\raggedright\arraybackslash}p{0.30\linewidth}
>{\raggedright\arraybackslash}p{0.24\linewidth}
>{\raggedright\arraybackslash}X@{}}
\toprule
\textbf{Family} & \textbf{Relation} & \textbf{Pair type} & \textbf{Temporal state} \\
\midrule
\multirow{4}{*}{\textbf{Physical}}
& \texttt{contact} & \texttt{ee--obj}, \texttt{obj--obj} & -- \\
& \texttt{grasp} & \texttt{ee--obj} & -- \\
& \texttt{support} & \texttt{obj--obj} & -- \\
& \texttt{contain} & \texttt{obj--obj} & -- \\
\midrule
\multirow{2}{*}{\textbf{Spatial}}
& \texttt{planar-distance} & \texttt{ee--obj}, \texttt{obj--obj} & change \\
& \texttt{height-offset} & \texttt{ee--obj}, \texttt{obj--obj} & change \\
\midrule
\multirow{4}{*}{\textbf{Affordance}}
& \texttt{contact-compatibility} & \texttt{ee--obj}, \texttt{obj--obj} & change \\
& \texttt{grasp-compatibility} & \texttt{ee--obj} & change \\
& \texttt{support-compatibility} & \texttt{obj--obj} & change \\
& \texttt{contain-compatibility} & \texttt{obj--obj} & change \\
\bottomrule
\end{tabularx}
\end{table}

\paragraph{Physical-state relations.}
The \texttt{contact} label is obtained by thresholding the simulator contact
force, and \texttt{grasp} is given by the simulator grasp predicate. These
relations take labels in \(\{\texttt{holds},\texttt{not-holds}\}\). For
object--object pairs, \texttt{support} is determined from contact and the
vertical component of the pairwise force, while \texttt{contain} tests whether
the containee's mined key point lies inside the container's entry volume.
Because support and containment are directed, their labels are
\(\{\texttt{not-holds},\texttt{src-holds},\texttt{dst-holds}\}\), where the
endpoints follow a fixed ordering. The physical predicates are evaluated
independently, so several may hold for the same pair. We do not assign a
separate temporal label to physical relations, since their transitions are
already represented by consecutive absolute states.

\paragraph{Spatial relations.}
For an end-effector--object pair, \texttt{planar-distance} is the horizontal
distance between the TCP and the object centroid, and \texttt{height-offset}
is their signed vertical displacement. For an object--object pair, we use the
corresponding mined interaction-surface anchors when available. This avoids
measuring support or containment geometry from body-frame origins that may lie
far from the relevant surfaces. The planar-distance labels are
\(\{\texttt{very-near},\texttt{near},\texttt{medium},\texttt{far},
\texttt{very-far}\}\), and the height-offset labels are
\(\{\texttt{far-below},\texttt{below},\texttt{level},\texttt{above},
\texttt{far-above}\}\). Their bin edges are derived from successful-rollout
statistics and fixed before training. End-effector--object and
object--object measurements are calibrated separately because they occupy
different spatial ranges. A spatial relation without a mined calibration is
not emitted.

\paragraph{Affordance relations.}
Successful interactions define a set of object-local components for each
affordance. Grasp components contain an interaction anchor, approach direction,
and gripper width; contact components contain surface anchors and normals;
support components describe a support surface and the corresponding object
bottom; and containment components describe an entry region and the
containee's key axis. Multiple components preserve distinct successful
interaction modes.

For each admissible pair, we select the applicable component or component pair
and compute a normalized incompatibility score
\begin{equation}
d_{ij,t}^{\rho}
=
\frac{1}{|\mathcal M_{ij,t}^{\rho}|}
\sum_{m\in\mathcal M_{ij,t}^{\rho}}
\min\!\left(\frac{\epsilon_{ij,t}^{m}}{\alpha_m},1\right),
\qquad d_{ij,t}^{\rho}\in[0,1],
\label{eq:affordance_score}
\end{equation}
where \(\epsilon_{ij,t}^{m}\) is an available relation-specific mismatch and
\(\alpha_m\) is its normalization scale. The mismatches comprise position,
orientation, and gripper width for grasp; anchor position and surface
orientation for contact; planar alignment, vertical gap, and surface
orientation for support; and radial offset, axial offset, and entry alignment
for containment. Missing optional components are excluded from the average.
A lower score denotes a configuration closer to a demonstrated successful
interaction.

Fine-grained compatibility is exposed after the pair enters the
\texttt{very-near} or \texttt{near} planar-distance bins. The absolute state is
\begin{equation}
\sigma_{ij,t}^{\rho}
=
\begin{cases}
\texttt{unobserved},
& \text{the pair is not near},\\
\texttt{match},
& \text{the pair is near and } d_{ij,t}^{\rho}<\frac{1}{3},\\
\texttt{partial-match},
& \text{the pair is near and }
  \frac{1}{3}\le d_{ij,t}^{\rho}<\frac{2}{3},\\
\texttt{poor-match},
& \text{the pair is near and } d_{ij,t}^{\rho}\ge\frac{2}{3}.
\end{cases}
\label{eq:affordance_bins}
\end{equation}
Thus, \texttt{unobserved} denotes an admissible pair outside the fine-grained
evaluation range, whereas \texttt{poor-match} denotes an evaluated but
incompatible configuration. Whenever the score is measurable, we retain its
continuous value for temporal labeling even if the absolute state is
\texttt{unobserved}. For object pairs already connected by a physical relation
at episode initialization, we omit their compatibility relations so
that the affordance set represents interaction configurations that can be
acted upon rather than static scene layout.

\paragraph{Temporal-change relations.}
For each spatial or affordance relation with a contiguous history of raw
measurements, we compute
\begin{equation}
\Delta u_{ij,t}^{\rho}
=
u_{ij,t}^{\rho}-u_{ij,t-K}^{\rho},
\label{eq:temporal_difference}
\end{equation}
where \(u_{ij,t}^{\rho}\) is the continuous spatial measurement or
incompatibility score underlying the absolute label. Negative changes in
planar distance denote approach, negative changes in height offset denote
downward relative motion, and negative changes in incompatibility denote
improving affordance alignment.

Let \(m_{\rho}>0\) be the robust change scale mined for relation \(\rho\).
We discretize the signed difference as
\begin{equation}
\delta_{ij,t}^{\rho}
=
\begin{cases}
\texttt{decrease-fast},
& \Delta u_{ij,t}^{\rho}< -3m_{\rho}/5,\\
\texttt{decrease-slow},
& -3m_{\rho}/5\le\Delta u_{ij,t}^{\rho}< -m_{\rho}/12,\\
\texttt{stable},
& |\Delta u_{ij,t}^{\rho}|\le m_{\rho}/12,\\
\texttt{increase-slow},
& m_{\rho}/12<\Delta u_{ij,t}^{\rho}\le 3m_{\rho}/5,\\
\texttt{increase-fast},
& \Delta u_{ij,t}^{\rho}>3m_{\rho}/5.
\end{cases}
\label{eq:temporal_bins}
\end{equation}
The change scales are calibrated separately for each affordance relation and
for each spatial endpoint type. A temporal label is produced only after
\(K+1\) consecutive valid measurements. If a relation is omitted or its raw
measurement is unavailable, its history is cleared and
\(\delta_{ij,t}^{\rho}=\varnothing\) until a new window is accumulated. The
affordance near gate does not by itself clear this history: continuous scores
can therefore reveal movement toward or away from an interaction-ready
configuration before the coarse absolute label changes.
